% =====================================================================
\chapter{Εισαγωγή}
\label{ch:intro}
% =====================================================================

\section{Κίνητρο}
\label{sec:motivation}

Το \en{Universal Serial Bus} (\en{USB}) είναι σήμερα η κυρίαρχη διεπαφή σύνδεσης
περιφερειακών σε προσωπικούς υπολογιστές, φορητές συσκευές και ενσωματωμένα
συστήματα. Η επιτυχία του οφείλεται σε μεγάλο βαθμό στη λογική του
«\en{plug and play}»: ο χρήστης συνδέει μια συσκευή και εκείνη λειτουργεί, χωρίς
καμία ρύθμιση. Από τη σκοπιά της ασφάλειας, ωστόσο, η ίδια ιδιότητα σημαίνει ότι
\emph{το λειτουργικό σύστημα εμπιστεύεται αυτόματα ό,τι δηλώνει η συσκευή για τον
εαυτό της}~\cite{tian2018sok}.

Μια συσκευή \en{USB} περιγράφει τον εαυτό της μέσω \emph{περιγραφέων}
(\en{descriptors}): δηλώνει κατασκευαστή, προϊόν, σειριακό αριθμό και ένα σύνολο
\emph{διεπαφών} (\en{interfaces}), καθεμία από τις οποίες ανήκει σε μια κλάση,
όπως «μαζική αποθήκευση» ή «συσκευή διεπαφής ανθρώπου» (\en{HID}, π.χ.
πληκτρολόγιο). Τίποτα στο πρωτόκολλο δεν επαληθεύει αυτές τις δηλώσεις. Το 2014,
οι \en{Nohl} και \en{Lell} έδειξαν ότι το υλικολογισμικό (\en{firmware}) κοινών
ελεγκτών στικακιών μνήμης μπορεί να επαναπρογραμματιστεί ώστε η συσκευή να
εμφανίζεται \emph{και} ως πληκτρολόγιο, ονομάζοντας την επίθεση
\en{BadUSB}~\cite{nohl2014badusb}. Εμπορικά διαθέσιμα εργαλεία, όπως το
\en{USB Rubber Ducky}~\cite{hak5ducky} και το καλώδιο \en{O.MG}~\cite{omgcable},
καθιστούν σήμερα τέτοιες επιθέσεις προσιτές σε οποιονδήποτε.

Η απειλή δεν είναι θεωρητική. Η μελέτη των \en{Tischer et~al.}~\cite{tischer2016usbdrop}
έδειξε ότι από 297 στικάκια που αφέθηκαν σε πανεπιστημιούπολη, το 45\% συνδέθηκαν
σε υπολογιστή και ανοίχτηκαν αρχεία τους, ενώ το \en{Stuxnet} χρησιμοποίησε
αφαιρούμενα μέσα \en{USB} για να διασχίσει απομονωμένα δίκτυα~\cite{falliere2011stuxnet}.
Μια συσκευή έγχυσης πληκτρολογήσεων μπορεί να εκτελέσει ολόκληρο σενάριο εντολών
σε λιγότερο από ένα δευτερόλεπτο, πολύ πριν ο χρήστης αντιληφθεί ότι κάτι συμβαίνει.

Στο \en{Linux}, ο πυρήνας προσδένει οδηγό σε μια νέα συσκευή μόλις ολοκληρωθεί η
απαρίθμησή της. Ένα ψεύτικο πληκτρολόγιο αποκτά συνεπώς πρόσβαση στη συνεδρία του
χρήστη χωρίς καμία ερώτηση. Ο πυρήνας ωστόσο προσφέρει, ήδη από το 2007, έναν
μηχανισμό \emph{εξουσιοδότησης} συσκευών μέσω του \en{sysfs}~\cite{kernel-usb-authorization}:
αν η τιμή \code{authorized\_default} ενός ριζικού διανομέα (\en{root hub}) τεθεί
σε 0, κάθε νέα συσκευή κάτω από αυτόν εμφανίζεται με \code{authorized=0}, δηλαδή
χωρίς οδηγό και χωρίς δυνατότητα να επικοινωνήσει με το σύστημα. Ο μηχανισμός
αυτός είναι το θεμέλιο της παρούσας εργασίας.

\section{Διατύπωση του προβλήματος}
\label{sec:problem}

Το ερώτημα που απαντά η εργασία είναι:

\begin{notebox}[Ερευνητικό ερώτημα]
Πώς μπορεί ένας σταθμός εργασίας \en{Linux} να αποφασίζει \emph{πριν} από την
ενεργοποίηση μιας νέας συσκευής \en{USB} αν αυτή πρέπει να γίνει αποδεκτή, έτσι
ώστε (α) καμία συσκευή να μη λειτουργεί χωρίς ρητή απόφαση, (β) η απόφαση να
στηρίζεται σε τεκμηριωμένη ανάλυση των ισχυρισμών και της συμπεριφοράς της
συσκευής, και (γ) ο χρήστης να μην κινδυνεύει ποτέ να κλειδωθεί έξω από το
σύστημά του;
\end{notebox}

Οι υπάρχουσες λύσεις καλύπτουν μέρος του προβλήματος. Το \en{USBGuard}~\cite{usbguard}
εφαρμόζει πολιτικές λευκής λίστας, αλλά δεν εξετάζει αν οι ισχυρισμοί μιας
άγνωστης συσκευής είναι αντιφατικοί, ούτε παρατηρεί τη συμπεριφορά της. Ερευνητικά
πρωτότυπα όπως τα \en{GoodUSB}~\cite{tian2015goodusb}, \en{USBFILTER}~\cite{tian2016usbfilter}
και \en{Cinch}~\cite{angel2016cinch} απαιτούν τροποποιήσεις στον πυρήνα ή
εικονικοποίηση, γεγονός που δυσχεραίνει την υιοθέτησή τους σε καθημερινούς
σταθμούς εργασίας. Το κενό που στοχεύει η εργασία είναι ένα εργαλείο που
\emph{τρέχει σε αμετάβλητο πυρήνα}, \emph{συνδυάζει} ελέγχους ταυτότητας,
συνέπειας, συμπεριφοράς και αποθήκευσης, και είναι \emph{ασφαλές στην αποτυχία}
(\en{fail-safe}) ως προς την πρόσβαση του νόμιμου χρήστη.

\section{Στόχοι της εργασίας}
\label{sec:goals}

Οι στόχοι της εργασίας είναι:

\begin{enumerate}[label=\textbf{Σ\arabic*.}]
  \item \textbf{Προεπιλεγμένη άρνηση.} Κάθε νέα συσκευή \en{USB} να παραμένει
        απενεργοποιημένη έως ότου ληφθεί απόφαση από άνθρωπο, από ρητά
        απομνημονευμένη εμπιστοσύνη ή από ρητή εξαίρεση ασφαλείας.
  \item \textbf{Τεκμηριωμένη απόφαση.} Πριν από την ερώτηση, να παρουσιάζεται στον
        χρήστη αναφορά με ευρήματα κλιμακούμενης σοβαρότητας
        (\sevC, \sevW, \sevN).
  \item \textbf{Ανίχνευση αντιφάσεων.} Να εντοπίζονται συνδυασμοί λειτουργιών
        χαρακτηριστικοί επιθέσεων, όπως «αποθήκευση + πληκτρολόγιο».
  \item \textbf{Παρατήρηση συμπεριφοράς.} Να ανιχνεύεται συσκευή που
        πληκτρολογεί μόνη της, χωρίς οι πληκτρολογήσεις της να φτάνουν στη
        συνεδρία του χρήστη.
  \item \textbf{Μνήμη.} Να διατηρείται ιστορικό ώστε να εντοπίζεται συσκευή
        γνωστής ταυτότητας που άλλαξε περιγραφείς.
  \item \textbf{Ελάχιστο προνόμιο.} Η ανάλυση μη έμπιστων δεδομένων να γίνεται
        σε διεργασία χωρίς προνόμια διαχειριστή.
  \item \textbf{Αποφυγή κλειδώματος.} Η πύλη να επανανοίγει σε κάθε τερματισμό,
        κατάρρευση ή σήμα, και να υπάρχουν μηχανισμοί διάσωσης.
  \item \textbf{Επαληθευσιμότητα.} Κάθε ισχυρισμός του συστήματος να
        τεκμηριώνεται με αυτοματοποιημένους ελέγχους, εκτελέσιμους χωρίς υλικό.
\end{enumerate}

\section{Συνεισφορά}
\label{sec:contributions}

Η συνεισφορά της εργασίας συνοψίζεται ως εξής:

\begin{itemize}
  \item \textbf{Σχεδίαση μιας πύλης τεσσάρων σταδίων} (ταυτότητα, συνέπεια,
        συμπεριφορά, αποθήκευση) πάνω στη διεπαφή εξουσιοδότησης του πυρήνα,
        χωρίς τροποποίηση του πυρήνα (Κεφάλαιο~\ref{ch:design}).
  \item \textbf{Ένα σύνολο 21 κανόνων συνέπειας} για περιγραφείς \en{USB},
        8 κανόνων συμπεριφοράς, 8 κανόνων αποθήκευσης, 3 κανόνων ιστορικού
        και 4 κανόνων για κάρτες μνήμης, με ρητό μοντέλο σοβαρότητας
        (Παράρτημα~\ref{app:rules}).
  \item \textbf{Μια αρχιτεκτονική διαχωρισμού προνομίων} όπου η προνομιούχος
        πύλη αντλεί το πεδίο δράσης της από την οπτική του πυρήνα και όχι από το
        αίτημα του αναλυτή, και όπου η «μόνιμη εμπιστοσύνη» εγγράφεται μόνο με
        αποτύπωμα που μέτρησε η ίδια η πύλη (§\ref{sec:design-privsep}).
  \item \textbf{Ένα μοντέλο κλιμακούμενης τριβής} στην αλληλεπίδραση με τον
        χρήστη: ένας διάλογος, δύο διάλογοι ή αντίστροφη μέτρηση 10 δευτερολέπτων,
        ανάλογα με τον κίνδυνο (§\ref{sec:design-friction}).
  \item \textbf{Μια υλοποίηση ανοικτού κώδικα} σε \en{Python} (περίπου 17\,700
        γραμμές κώδικα παραγωγής) με 1413 αυτοματοποιημένους ελέγχους,
        ελέγχους ιδιοτήτων, \en{fuzzing} και πλαίσιο εξομοίωσης συσκευών
        (Κεφάλαια~\ref{ch:implementation} και~\ref{ch:evaluation}).
  \item \textbf{Ένα τεκμηριωμένο αρχείο ποιότητας} 106 ελαττωμάτων με τη μέθοδο
        εντοπισμού και τη σοβαρότητά τους, ως εμπειρικά δεδομένα για την
        αποτελεσματικότητα διαφορετικών τεχνικών επαλήθευσης (§\ref{sec:eval-qa}).
\end{itemize}

\section{Δομή της εργασίας}
\label{sec:structure}

Η εργασία οργανώνεται ως εξής. Το Κεφάλαιο~\ref{ch:background} παρουσιάζει το
απαραίτητο θεωρητικό υπόβαθρο: την αρχιτεκτονική \en{USB}, τους περιγραφείς, τη
διαδικασία απαρίθμησης και τους μηχανισμούς του \en{Linux} που αξιοποιούνται
(\en{sysfs}, \en{udev}, \en{evdev}). Το Κεφάλαιο~\ref{ch:related} ορίζει το μοντέλο
απειλής και επισκοπεί τη σχετική βιβλιογραφία. Το Κεφάλαιο~\ref{ch:design}
περιγράφει τις απαιτήσεις και τη σχεδίαση του συστήματος, ενώ το
Κεφάλαιο~\ref{ch:implementation} την υλοποίησή του. Το Κεφάλαιο~\ref{ch:evaluation}
παρουσιάζει την πειραματική αξιολόγηση με παραδείγματα σεναρίων επίθεσης. Το
Κεφάλαιο~\ref{ch:discussion} συζητά περιορισμούς και συγκρίνει με υπάρχουσες
λύσεις, και το Κεφάλαιο~\ref{ch:conclusions} συνοψίζει τα συμπεράσματα και
προτείνει μελλοντική εργασία. Τα Παραρτήματα περιέχουν οδηγίες εγκατάστασης, τον
πλήρη κατάλογο κανόνων και παραδείγματα εξόδου.
